\paragraph{Theorem F.2.24 (Proved: every surviving twist in the unbounded family).}
Use the fixed Kuhn cusp and the affine translation clutching
convention of F.2.12.  In the entire rational family F.2.16,
allow every stable lattice, every free affine finite lift,
and every pair of actual finite multiplicities.
The data with $\pi_1=1$ and $H_2=H_3=0$ are precisely the
following, up to the integral linear classification already
proved and the translation identifications described below:
\[
\begin{gathered}
 v\in\{0,1\},\qquad
 \text{the unique eligible integral class of F.2.23},\\
 (m_1,m_2)=(5,4),\qquad
 20\kappa+4r_1+5r_2=\pm1.
\end{gathered}
 \tag{F.T1}
\]
Here $r_i=\ell(v_i)$ is the invariant affine numerator, and
$\kappa\in\mathbb Z$ is the invariant coordinate of the
cusp translation clutching, with signs fixed by the displayed
van Kampen relation below.  The remaining fixed-direction
components of the twists impose no condition.
Every manifold in (F.T1) is oriented diffeomorphic to its
corresponding known survivor and hence to
$3\mathbb{HP}^2\#(-\mathbb{HP}^2)$.
This completes the survivor classification for this entire
unbounded rational family; it does not classify the other
rational points of F.2.20 or the other spectral loci.

\begin{proof}
F.2.19 excludes all other rational parameters, and F.2.23
excludes all other integral linear classes.
We therefore work with either known integral pair, in the
marking of (F.U6).  Its primitive invariant homology
covector is $\ell(x)=x_0$.
Each finite fixed lattice has an integral basis
$u_i,\lambda_i$, where $\ell(u_i)=1$,
$\ell(\lambda_i)=0$, and $\lambda_i$ is primitive in
$\ker\ell$.
We may write the affine numerator as
$v_i=r_i u_i+s_i\lambda_i$.
This covers arbitrary finite affine lifts: translating the
origin makes their translation vector invariant by averaging,
and its $m_i$th iterate is then the integral vector $v_i$.
The resulting boundary translation change is included in the
specified affine clutching.

For both linear pairs, the normal closure of the cusp angular
lattice in the fibre lattice is exactly $\ker\ell$, by
F.2.5 and F.2.18.  Van Kampen therefore leaves the abelian
presentation
\[
 m_1a=r_1h,\qquad m_2b=r_2h,\qquad
 a+b+\kappa h=0.
 \tag{F.T2}
\]
This also defines the sign of $\kappa$.
Its determinant is, up to sign,
\[
 p=m_1m_2\kappa+m_2r_1+m_1r_2.
 \tag{F.T3}
\]
Thus simple connectivity is exactly $p=\pm1$.
It forces $\gcd(m_1,m_2)=1$ and
$\gcd(r_i,m_i)=1$ by reduction modulo each multiplicity.

We first prove that a surviving datum has the minimal
multiplicities.  Set $n_1=5$, $n_2=4$ and $m_i=n_i k_i$.
In the splitting $\Lambda=\mathbb Zu_i\oplus\ker\ell$,
write the finite linear map as $1\oplus D_i$.
The reduced finite fibre is a mapping torus of $T^5$ over
the circle $m_i\ell$; its fibre is connected because
$r_i$ is prime to $m_i$.
Its vertical monodromy is a power of $D_i$ prime to $n_i$,
possibly followed by a translation, which changes no
integral Wang group.

The primitive fixed vertical loop $\lambda_i$ gives the
only additional free two-cycle beyond vertical cycles.
Consequently an invariant integral two-form $\omega$ on
the covering torus descends integrally exactly when
\[
 \frac{\omega(v_i,\lambda_i)}{m_i}
       =\frac{r_i\omega(u_i,\lambda_i)}{m_i}
       \quad\text{is an integer}.
 \tag{F.T4}
\]
For precision, the torus made from the affine generator
and $\lambda_i$ projects with degree $r_i$ to the mapping-
torus base, whereas the ordinary $u_i$ suspension projects
with degree $m_i$.  A Bezout combination is a primitive
suspension.  Its integrality condition is equivalent to
(F.T4), because the ordinary covering periods are already
integral and $\gcd(r_i,m_i)=1$.
The Wang sequence shows that these vertical cycles and one
primitive suspension generate the free second homology;
there is no omitted period condition.

The integral invariant two-form computations in F.2.6 and
F.2.18 give
\[
 \gcd_{\omega}\omega(u_i,\lambda_i)=n_i.
 \tag{F.T5}
\]
They are also independently printed for both pairs by
\texttt{fourfolds\_family\_twists.py}.
Hence the covering-image cokernel in degree two at the
$i$th finite point is exactly $\mathbb Z/k_i$.
The global invariant group of $\bigwedge^2\mathcal V$
and its integral first parabolic group both vanish.
The skyscraper cokernel sequence therefore gives
\[
 H^1(B,R^2g_*\mathbb Z)
       =\mathbb Z/k_1\oplus\mathbb Z/k_2.
 \tag{F.T6}
\]
Point-supported torsion kernels do not change this degree-one
base cohomology, and the cusp specialization is onto.
The term (F.T6), in bidegree $(1,2)$, survives every Leray
differential on the two-dimensional base.
All rational terms of total degree three vanish, independently
of the multiplicities and twists, by the same exterior
local-system table as F.2.6.  Thus $H^3$ is finite.
If either $k_i>1$, its surviving nonzero filtration quotient
makes $H^3$ nonzero; universal coefficients then give
nonzero torsion in $H_2$.
This proves $m_1=5$, $m_2=4$ for every survivor.

We now identify all twists with these multiplicities.
The following integral fact permits removal of each $s_i$:
\[
 \operatorname{Norm}_{n_i}:
 \ker\ell\longrightarrow\mathbb Z\lambda_i
 \quad\text{is onto}.
 \tag{F.T7}
\]
Here the norm is $I+D_i+\cdots+D_i^{n_i-1}$.
One can read (F.T7) directly from (F.T5).
Let $\varphi_i$ generate the primitive invariant covector
lattice on $\ker\ell$.
The splitting by $u_i$ shows that the gcd in (F.T5) is
$|\varphi_i(\lambda_i)|$: its forms are
$\ell\wedge\varphi_i$ plus invariant vertical forms.
Thus that value is $n_i$.
The norm has the form
$x\mapsto d_i\varphi_i(x)\lambda_i$ for an integer
$d_i$; evaluating $\varphi_i$ gives
$|d_i|n_i=n_i$, hence $|d_i|=1$.
This proves surjectivity.  The saved program additionally
gives integral preimages $q_i\in\ker\ell$ with
$\operatorname{Norm}_{n_i}q_i=\lambda_i$.

Put
\[
 w_i=-\frac{s_i}{n_i}\sum_{j=0}^{n_i-1}jM_i^j q_i.
\]
Then
\[
 (I-M_i)w_i=-s_i\lambda_i/n_i+s_iq_i.
 \tag{F.T8}
\]
Translation by $w_i$ conjugates the finite affine action
with numerator $r_i u_i+s_i\lambda_i$ to that with
numerator $r_i u_i$: the remaining $s_iq_i$ is integral.
It gives a diffeomorphism of the entire finite cap.
In the standard linear boundary marking, its translation
section is
$w_i+\theta s_i\lambda_i/n_i$.
Its endpoint mismatch with the linear deck map is exactly
$s_iq_i$ by (F.T8), so its integral boundary residue has
$\ell$-coordinate zero.

We explain why these local changes extend simultaneously,
including the prescribed cusp marking.
Let $P$ be the complementary pair of pants and
$\mathcal T\to P$ its flat torus group bundle.
A section on the boundary extends over the relative
one-skeleton because the torus is connected.
The only remaining obstruction is the two-cell obstruction
with values in the fibre fundamental-group local system:
\[
 H^2(P,\partial P;\Lambda)
        \simeq H_0(P;\Lambda)=\Lambda_\Gamma.
 \tag{F.T9}
\]
The equality follows from the oriented relative cellular
complex, equivalently integral Poincare--Lefschetz duality
with local coefficients.  Its boundary residues map to
the indicated global coinvariants.
For both linear pairs the horizontal monodromy relation
matrix has five unit Smith factors, so
$\Lambda_\Gamma=\mathbb Z$, detected exactly by $\ell$.
This is independently checked in the saved program.
Thus the residues $s_iq_i$, with the identity section at the
cusp, have zero obstruction and extend over $P$.
A section may be smoothed relative to its boundary collars;
translation by it is a fibre-preserving diffeomorphism.
Together with the finite-cap maps it removes all $s_i$
without changing the cusp.
The same argument removes all components of translation
clutching whose $\ell$-coordinate is zero.

It remains to handle the integer numerators and the remaining
cusp coordinate.  If $p=-1$, the congruences from (F.T3)
give
\[
 r_1=1+5a,\qquad r_2=-1+4b,
 \qquad a+b+\kappa=0.
 \tag{F.T10}
\]
Adding $5a u_1$ or $4b u_2$ changes no finite affine action;
in its linear boundary marking it changes the translation
residue by $a u_1$ or $b u_2$.
The difference of the cusp marking has residue of
$\ell$-coordinate $\kappa$, by (F.T2).
Their total obstruction in (F.T9) is therefore
$a+b+\kappa=0$.
The same extension argument identifies the glued manifold
with the original twists $(u_1,-u_2)$ and zero cusp
clutching.
If $p=1$, the identical calculation starts from
$(-u_1,u_2)$.

Finally simultaneous fibre inversion identifies these two
sign choices by an orientation-preserving diffeomorphism.
It conjugates every affine translation to its negative.
It extends over the Kuhn cusp: the lattice involution
$(x,h)\mapsto(-x,h)$ preserves its toric fan, since a
simplex given by a chain in a unit cube is taken to the
chain with reversed coordinate order in the opposite cube.
It normalizes the cusp lattice action and therefore descends
to its quotient, while fixing the base parameter.  The resulting toric map is holomorphic,
and in the regular real six-torus inversion has determinant
$(-1)^6=1$.  All translation maps also preserve orientation.
Thus every datum in (F.T1) gives the stated oriented smooth
type and has $H_2=H_3=0$, proving sufficiency as well as
necessity.
\end{proof}
